\documentclass[10pt,a4paper]{amsart}
\usepackage{amsmath,amssymb,amsthm,graphicx}
% For commutative diagrams we use Xy-pic
\usepackage[all]{xy}

% We use 2cell for 2-commutative diagrams.
\xyoption{2cell}
% Theorem environments.
%
\theoremstyle{plain}
\newtheorem{theorem}[subsection]{Theorem}
\newtheorem{proposition}[subsection]{Proposition}
\newtheorem{lemma}[subsection]{Lemma}

\theoremstyle{definition}
\newtheorem{definition}[subsection]{Definition}
\newtheorem{example}[subsection]{Example}
\newtheorem{exercise}[subsection]{Exercise}
\newtheorem{situation}[subsection]{Situation}

\theoremstyle{remark}
\newtheorem{remark}[subsection]{Remark}
\newtheorem{remarks}[subsection]{Remarks}

\numberwithin{equation}{subsection}

% Macros
%
\def\lim{\mathop{\mathrm{lim}}\nolimits}
\def\colim{\mathop{\mathrm{colim}}\nolimits}
\def\Spec{\mathop{\mathrm{Spec}}}
\def\Hom{\mathop{\mathrm{Hom}}\nolimits}
\def\Ext{\mathop{\mathrm{Ext}}\nolimits}
\def\SheafHom{\mathop{\mathcal{H}\!\mathit{om}}\nolimits}
\def\SheafExt{\mathop{\mathcal{E}\!\mathit{xt}}\nolimits}
\def\Sch{\mathit{Sch}}
\def\Mor{\mathop{\mathrm{Mor}}\nolimits}
\def\Ob{\mathop{\mathrm{Ob}}\nolimits}
\def\Sh{\mathop{\mathit{Sh}}\nolimits}
\def\NL{\mathop{N\!L}\nolimits}
\def\CH{\mathop{\mathrm{CH}}\nolimits}
\def\proetale{{pro\text{-}\acute{e}tale}}
\def\etale{{\acute{e}tale}}
\def\QCoh{\mathit{QCoh}}
\def\Ker{\mathop{\mathrm{Ker}}}
\def\Im{\mathop{\mathrm{Im}}}
\def\Coker{\mathop{\mathrm{Coker}}}
\def\Coim{\mathop{\mathrm{Coim}}}

% Boxtimes
%
\DeclareMathSymbol{\boxtimes}{\mathbin}{AMSa}{"02}

%
% Macros for moduli stacks/spaces
%
\def\QCohstack{\mathcal{QC}\!\mathit{oh}}
\def\Cohstack{\mathcal{C}\!\mathit{oh}}
\def\Spacesstack{\mathcal{S}\!\mathit{paces}}
\def\Quotfunctor{\mathrm{Quot}}
\def\Hilbfunctor{\mathrm{Hilb}}
\def\Curvesstack{\mathcal{C}\!\mathit{urves}}
\def\Polarizedstack{\mathcal{P}\!\mathit{olarized}}
\def\Complexesstack{\mathcal{C}\!\mathit{omplexes}}
% \Pic is the operator that assigns to X its picard group, usage \Pic(X)
% \Picardstack_{X/B} denotes the Picard stack of X over B
% \Picardfunctor_{X/B} denotes the Picard functor of X over B
\def\Pic{\mathop{\mathrm{Pic}}\nolimits}
\def\Picardstack{\mathcal{P}\!\mathit{ic}}
\def\Picardfunctor{\mathrm{Pic}}
\def\Deformationcategory{\mathcal{D}\!\mathit{ef}}
\usepackage{cleveref}
\crefname{originallemma}{original lemma}{original lemmas}
\Crefname{originallemma}{Original lemma}{Original lemmas}
\makeatletter
\newlabel{lemma-det-exact-sequences}{{02PA}{source}{}{sourcepointer-1}{}}
\newlabel{lemma-det-exact-sequences@cref}{{[originallemma][0][]02PA}{[1][0][]source}{}{}{}}
\makeatother
\crefname{originallemma}{original lemma}{original lemmas}
\Crefname{originallemma}{Original lemma}{Original lemmas}
\makeatletter
\newlabel{algebra-lemma-pushdown-module}{{02M0}{source}{}{sourcepointer-2}{}}
\newlabel{algebra-lemma-pushdown-module@cref}{{[originallemma][0][]02M0}{[1][0][]source}{}{}{}}
\makeatother
\title{Canonical functorial isomorphism between the determinants of a finite length module over a local ring and over a local quotient-ring extension with the same residue field}
\author{The Stacks Project Authors}
\date{stacks-project:chow.tex:local-snapshot-20261012; GNU FDL 1.2 or later}
\begin{document}
\maketitle
\noindent\textit{Editorial note.} All delivered mathematical text is the author's own, copied from the pinned Stacks Project chapter file chow.tex; no statement or proof step was rewritten, reordered or removed.\par\medskip
\noindent\textit{Editorial note.} Preamble provenance: the pinned chapter file opens with an input of the shared preamble and carries no preamble text of its own, so the project's single Stacks-wide preamble staged at sources/stacks-chow/preamble.tex is excerpted verbatim as two spans, lines 14-18 (Xy-pic setup) and lines 154-217 (the native theorem declarations and mathematical macros). Nothing was added, removed or redefined; the wrapper class and the amsmath/amssymb/amsthm/graphicx base packages are supplied by the assembler.\par\medskip
\noindent\textit{Editorial note.} The counted statement and proof contain no cross reference, no citation and no included file, so the delivered body is self contained.\par\medskip
\noindent\textit{Editorial note.} Canonical location of the counted proposition: Stacks tag 02PC (sources/stacks-chow/pinned-original-source.tex line 16434).\par\medskip
\noindent\textit{Allowed standard prerequisites (editorial).}\begin{itemize}
\item Length of a finite length module is preserved under restriction along a local ring homomorphism. Algebra, Lemma labelled algebra-lemma-pushdown-module (Stacks tag 02M0), invoked at proof lines 16451-16452. Supplies that M\_\{R'\} has the same finite length l as M, which is what makes the two determinant modules one dimensional over the same residue field.
\item The determinant of a finite length module and the isomorphisms gamma for short exact sequences. Same chapter, Lemma labelled lemma-det-exact-sequences (source line 16203), invoked in the statement and at proof line 16462. Supplies the determinant construction and the compatibility assertion whose verification is claimed in the proof; the symbol-wise definition of the isomorphism is inside the counted span.
\end{itemize}\medskip
\ifdefined\hypertarget\hypertarget{sourcepointer-1}{}\fi
\noindent\textit{Source reference.} Original lemma 02PA: sources/stacks-chow/pinned-original-source.tex, source lines 16203--16203. Same-chapter pointer into the author environment lemma at source line 16203; invoked at proof line 16462 of the counted proof. It is retained only as a navigation pointer; every load-bearing step is inside the delivered span.\par\medskip
\ifdefined\hypertarget\hypertarget{sourcepointer-2}{}\fi
\noindent\textit{Source reference.} Original lemma 02M0: sources/stacks-chow/tags.txt, source lines 3427--3427. Cross-chapter standard prerequisite located in the project tag list staged in the local snapshot; invoked at proof line 16452 of the counted proof. It is retained only as a navigation pointer; every load-bearing step is inside the delivered span.\par\medskip
\begin{lemma}
\label{lemma-determinant-quotient-ring}
Let $(R', \mathfrak m') \to (R, \mathfrak m)$ be a local ring
homomorphism which induces an isomorphism on residue fields $\kappa$.
Then for every finite length $R$-module the restriction $M_{R'}$
is a finite length $R'$-module and there is a canonical isomorphism
$$
\det\nolimits_{R, \kappa}(M)
\longrightarrow
\det\nolimits_{R', \kappa}(M_{R'})
$$
This isomorphism is functorial in $M$ and compatible with the
isomorphisms $\gamma_{K \to L \to M}$ of Lemma \ref{lemma-det-exact-sequences}
defined for $\det_{R, \kappa}$ and $\det_{R', \kappa}$.
\end{lemma}

\begin{proof}
If the length of $M$ as an $R$-module is $l$, then the length
of $M$ as an $R'$-module (i.e., $M_{R'}$) is $l$ as well, see
Algebra, Lemma \ref{algebra-lemma-pushdown-module}.
Note that an admissible sequence $x_1, \ldots, x_l$ of $M$
over $R$ is an admissible sequence of $M$ over $R'$ as $\mathfrak m'$
maps into $\mathfrak m$.
The isomorphism is obtained by mapping the symbol
$[x_1, \ldots, x_l] \in \det\nolimits_{R, \kappa}(M)$
to the corresponding symbol
$[x_1, \ldots, x_l] \in \det\nolimits_{R', \kappa}(M)$.
It is immediate to verify that this is functorial for
isomorphisms and compatible with the isomorphisms
$\gamma$ of Lemma \ref{lemma-det-exact-sequences}.
\end{proof}

\end{document}
